% =====================================================================
%  IGARSS submission.
%
%  TEMPLATE: verify against the official IGARSS paper kit before
%  submitting. IGARSS has historically used the ICASSP-style spconf
%  kit (spconf.sty + IEEEbib.sty), not IEEEtran. This file is set up
%  for IEEEtran so it compiles anywhere as-is; if the kit mandates
%  spconf, make the four swaps marked "SPCONF" below.
%
%  PAGE LIMIT: IGARSS is normally 4 pages INCLUDING references.
%  This draft is close to that limit -- check the compiled length
%  before adding anything.
% =====================================================================

\documentclass[conference]{IEEEtran}
% SPCONF 1/4 -- replace the line above with:
%   \documentclass{article}
%   \usepackage{spconf}

\usepackage{amsmath, amssymb}
\usepackage{graphicx}
\usepackage{booktabs}
\usepackage[hidelinks]{hyperref}
\usepackage{cite}

\hypersetup{colorlinks=true,linkcolor=blue!60!black,citecolor=blue!60!black,urlcolor=blue!60!black}

\title{A Land-Cover Data-Ceiling Diagnostic for Zero-Authentication
Urban Heat Island Modelling}

\author{
\IEEEauthorblockN{Abhiram Donepudi}
\IEEEauthorblockA{Westwood High School, Round Rock, TX\\
\texttt{abhiram.donepudi@gmail.com}}
}
% SPCONF 2/4 -- replace the \author{...} block above with:
%   \name{Abhiram Donepudi}
%   \address{Westwood High School, Round Rock, TX \\
%            \texttt{abhiram.donepudi@gmail.com}}

\begin{document}
\maketitle

\begin{abstract}
Whether a low urban heat island (UHI) transfer score reflects model
failure or a limit of the input data is normally indistinguishable. We
introduce a \emph{land-cover data-ceiling control} that separates the two
and is independent of any particular simulator: it requires a city
boundary, a land-cover raster and one satellite scene --- no model run, no
calibration, no meteorology. We define it as the multiple correlation
between MODIS land surface temperature (LST) and the land-cover channels a
model consumes, and show the simpler univariate alternative (impervious
fraction alone) is unsound, going negative in arid cities even where the
sound ceiling exceeds $R{=}{+}0.45$. Across 49 US cities a
built--vegetated thermal-gap heuristic predicts the ceiling (Spearman
$\rho{=}{+}0.69$), and extending to 28 cities on five continents under a
uniform city definition shows the arid inversion is climatic rather than
an artifact of an all-US sample. On a three-city transfer experiment,
Denver and Phoenix reach 29\% and 28\% of their multi-day ceilings while
Miami reaches 99\%.
\end{abstract}

\begin{IEEEkeywords}
Urban heat island, land surface temperature, MODIS, class separability,
model transferability
\end{IEEEkeywords}
% SPCONF 3/4 -- replace the IEEEkeywords environment above with:
%   \begin{keywords}
%   Urban heat island, land surface temperature, MODIS, class
%   separability, model transferability
%   \end{keywords}

\section{Introduction}

A recurring problem in UHI model evaluation is conflating two failure
modes when a model transfers to a new city: the model may genuinely fail
to capture the city's thermal structure, or the input land-cover data may
not separate thermal regimes well enough for \emph{any} land-cover-driven
model to succeed. This paper introduces a \emph{land-cover data-ceiling
control} that makes the distinction explicit, demonstrated on a
three-city transfer experiment (Miami, Phoenix, Denver) using
\textsc{austin-twin}, a physics-based slab energy-balance simulator of
Austin, TX driven entirely by zero-authentication open data
\cite{oke1982energetic, voogt2003thermal, masson2000physically}. Existing
UHI tools span two extremes: high-fidelity mesoscale models (WRF-SLUCM,
TEB, Urban Weather Generator \cite{bueno2013urban, kusaka2001simple})
require proprietary or authenticated inputs and hours of compute;
machine-learning approaches \cite{lyu2022integrated, hoang2025geospatial,
bhamjee2024detection} validate well historically but cannot predict an
intervention's effect outside the training distribution.
\textsc{austin-twin} occupies the middle: a physics-based 2D grid model
running in $10.8$\,ms per 48-hour simulation on a desktop CPU, with an
explicit counterfactual engine and fully open data. Against that
literature the data-ceiling diagnostic has no prior entry, and it is this
paper's primary contribution, independent of the specific simulator used
to generate the transfer scores.

The contribution is not a new physical formulation but three reusable
methodological tools, each independent of this simulator: (1) a
\emph{land-cover data-ceiling control} that must be defined as a multiple
correlation across all channels a model consumes --- the univariate
alternative is unsound, going negative in arid cities --- together with a
built--vegetated thermal-gap heuristic validated across 49 cities and an
NDVI channel screen shown to be a predictable, general rule rather than a
single-city result; (2) a decay-constrained two-stage calibration
quantifying the trade-off between MODIS pattern fidelity and intervention
spatial realism; and (3) an evapotranspiration ablation showing that a
standard linear ET term contributes nothing to spatial-pattern skill at
500\,m, regardless of formulation.

\section{Data and Methods}

\subsection{Data}
All inputs are retrieved anonymously without API keys. Austin's boundary
comes from OpenStreetMap via \texttt{osmnx}. Land cover uses ESA
WorldCover v200 (10\,m) \cite{zanaga2022esa}, aggregated to a 500\,m UTM
grid ($\sim$2{,}896 cells) into three channels: impervious, vegetation,
and water fraction. MODIS Aqua MYD11A1 v6.1 daily LST (1:30\,pm overpass)
is retrieved via Microsoft Planetary Computer \cite{microsoft_pc}. ERA5
hourly reanalysis is retrieved from Google's ARCO-ERA5 Zarr mirror
\cite{hersbach2020era5}. The same pipeline is applied to Phoenix, AZ;
Miami, FL; and Denver, CO for the transfer experiment, and MODIS NDVI
(MOD13A1, 500\,m) for the channel screen.

\subsection{Energy Balance Simulator}
Each cell evolves under
\begin{equation}
  \frac{dT_i}{dt} = Q_{\mathrm{solar}} - Q_{\mathrm{rad}}
                  - Q_{\mathrm{ET}} + Q_{\mathrm{diff}},
  \label{eq:slab}
\end{equation}
with per-class solar absorption, Newtonian radiative cooling, linear
evapotranspiration $Q_{\mathrm{ET}} = k_{\mathrm{ET}} f_{\mathrm{veg}}
\max(T_i - T_{\mathrm{air}}, 0)$, and a CFL-stable diffusion term. Water
cells use explicit mixed-layer heat capacity and a fixed latent-heat flux
rather than a single relaxation parameter.

\subsection{Calibration and Ceiling Control}
Five coefficients are tuned by differential evolution
\cite{storn1997differential} against three training MODIS days, with two
held out. The \emph{ship configuration} uses a two-stage procedure: fix
$D = 1.0\,\text{m}^2\text{s}^{-1}$ (matching the literature tree-cooling
half-distance \cite{ziter2019scale}), then refit the remaining four
coefficients under ARCO-ERA5 forcing. Ship coefficients transfer to
Phoenix, Miami, and Denver under three schemes: zero-shot; Scheme~B
(re-optimize $k_{\mathrm{ET}}$ only); Scheme~A (re-optimize all four
coefficients). A \emph{data-ceiling control} establishes the maximum
correlation any \emph{linear} model of the input channels can achieve. We
define it as the multiple correlation $R$ between LST and all three input
channels, not the univariate correlation with impervious fraction alone.
The univariate version is unsound: in arid cities it can be negative even
where the sound three-channel fit exceeds $R{=}{+}0.45$, because desert
vegetation runs hotter than built surfaces there
(Section~\ref{sec:results}). Because single-day ceilings are noisy, we
report multi-day means throughout (5--12 cloud-free dates per focus
city; median 6 across the 49-city sample).

Skill is reported as the spatial correlation and RMSE of mean-subtracted
fields. The slab formulation carries no absolute-temperature anchor, so
only the spatial pattern is meaningfully scored; RMSE computed on raw
fields would additionally carry a constant offset that the model is not
attempting to predict.

\section{Results}
\label{sec:results}

\subsection{Calibration and Uncertainty}
Table~\ref{tab:bootstrap_ci} reports spatial block-bootstrap 95\% CIs
($B{=}2{,}000$, $6{\times}6$-cell blocks). The ship configuration achieves
$r{=}{+}0.754$ [$+0.718$, $+0.785$], trading $0.10$ in $r$ against an
unconstrained single-stage fit ($r{=}{+}0.850$) for physically realistic
diffusion and intervention decay lengths (Section~\ref{sec:disc}). ERA5
forcing contributes $+0.007$ [$+0.004$, $+0.009$] over synthetic forcing ---
small but distinguishable from zero, with the benefit being mainly
methodological (absolute-temperature comparability without anomaly
normalization).

\begin{table}[ht]
  \centering
  \caption{Block-bootstrap 95\% CIs, held-out test $r$ and RMSE.}
  \label{tab:bootstrap_ci}
  \scriptsize
  \begin{tabular}{lcc}
    \toprule
    Configuration & Test $r$ [95\% CI] & RMSE \\
    \midrule
    Single-stage, synthetic forcing & $+0.850$ [$+0.823,+0.872$] & $1.09$\,$^{\circ}$C \\
    Two-stage, synthetic forcing & $+0.747$ [$+0.712,+0.778$] & $1.48$\,$^{\circ}$C \\
    \textbf{Ship (ARCO-ERA5)} & $\mathbf{+0.754}$ [$+0.718,+0.785$] & $\mathbf{1.40}$\,$^{\circ}$C \\
    \bottomrule
  \end{tabular}
\end{table}

\subsection{Multi-City Transfer and Ceiling Diagnostic}
Table~\ref{tab:transfer} presents the transfer matrix. Miami reaches 99\%
of its multi-day ceiling via Scheme~A ($r{=}{+}0.832$ vs.\
$+0.841\pm0.025$); Phoenix and Denver reach 28\% and 29\%
($+0.271\pm0.085$ and $+0.437\pm0.085$). Re-optimization degrades Phoenix
($+0.077\to+0.033$) but roughly doubles Denver ($+0.056\to+0.126$) ---
Denver's higher ceiling leaves headroom for calibration to exploit that
Phoenix's does not, illustrating that the diagnostic's practical use is
distinguishing fixable transfer failures from unfixable ones, not merely
flagging low scores.

\begin{table}[ht]
  \centering
  \caption{Multi-city transfer: held-out test $r$ (2 days). Ceiling $=$ multi-day mean $R$(LST; impervious, vegetation, water).}
  \label{tab:transfer}
  \scriptsize
  \setlength{\tabcolsep}{3pt}
  \begin{tabular}{lccc}
    \toprule
    Scheme & Denver & Phoenix & Miami \\
    \midrule
    Austin ship (ref.) & $+0.754$ & $+0.754$ & $+0.754$ \\
    Zero-shot        & $+0.056$ & $+0.077$ & $+0.724$ \\
    Scheme B, 3 days & $+0.063$ & $+0.071$ & $+0.776$ \\
    Scheme A, 3 days & $+0.126$ & $+0.033$ & $+0.832$ \\
    \midrule
    \textbf{Ceiling} & $\mathbf{+0.437\pm0.085}$ & $\mathbf{+0.271\pm0.085}$ & $\mathbf{+0.841\pm0.025}$ \\
    \bottomrule
  \end{tabular}
\end{table}

The built--vegetated LST gap (mean LST, impervious cells minus vegetated
cells) is the physical driver behind the ceiling, and validates on a
wider sample of 49 US cities (55 attempted, 6 dropped for cloud cover).
Against the sound multivariate ceiling, Spearman $\rho{=}{+}0.69$
($p{=}3\times10^{-8}$); a univariate proxy target gives a more flattering
$\rho{=}{+}0.84$, confirmed significantly higher by a paired bootstrap
($\Delta\rho{=}{+}0.15$, [$+0.03,+0.29$]) --- part of the univariate
figure's strength is a shared arid sign-flip, not gap predictiveness as
such. The heuristic generalizes across arid-exclusion subsets
($p{<}10^{-4}$) except the humid half ($n{=}24$, $\rho{=}{+}0.36$,
$p{=}0.08$), which we leave open: at this sample size, underpowering and a
genuine concentration in arid cities are not separable. In arid climates
the land-cover--temperature relationship \emph{inverts} rather than merely
weakens: natural desert vegetation runs hotter than irrigated, shaded
built surfaces, driving the univariate ceiling negative (Fig.~\ref{fig:inversion}).

\begin{table}[ht]
  \centering
  \caption{Cities with negative multi-day mean univariate ceiling, under the
  uniform 25\,km disc definition. Inversion appears on four continents.}
  \label{tab:arid}
  \scriptsize
  \begin{tabular}{llcc}
    \toprule
    City & Region & Ceiling & Gap ($^{\circ}$C) \\
    \midrule
    Amman & Levant & $-0.429$ & $-2.52$ \\
    Riyadh & Arabia & $-0.342$ & $-1.56$ \\
    Marrakesh & N.~Africa & $-0.156$ & $-1.18$ \\
    Antofagasta & S.~America & $-0.124$ & $-3.38$ \\
    Hermosillo & N.~America & $-0.116$ & $-0.99$ \\
    Alice Springs & Australia & $-0.096$ & $-3.11$ \\
    Cairo & N.~Africa & $-0.082$ & $+1.08$ \\
    Yuma & N.~America & $-0.067$ & $-0.48$ \\
    Mendoza & S.~America & $-0.032$ & $-0.57$ \\
    \bottomrule
  \end{tabular}
\end{table}

\begin{figure}[t]
  \centering
  \includegraphics[width=\columnwidth]{fig_inversion.pdf}
  \caption{The inversion as surface relief; height and colour both encode
  the plotted variable. The built core rises as a plateau in \emph{both}
  cities' impervious channel (left, range $0$--$1$), but the temperature
  surface rises with it in Austin ($38$--$48$\,$^{\circ}$C) and sinks into a basin
  in Riyadh ($53$--$57$\,$^{\circ}$C), where the surrounding desert is hotter than
  the city. Nothing in the land-cover input distinguishes the two cases,
  which is why a univariate land-cover ceiling reports that no model can
  work in the arid one. Single representative cloud-free dates; Riyadh
  under the 25\,km disc definition.}
  \label{fig:inversion}
\end{figure}

\subsection{Is Inversion Climatic or Continental?}
The 49-city sample is entirely American, so ``arid'' is confounded with
``American'' and inversion could reflect a continental artifact. We
therefore repeat the diagnostic on 28 cities across five continents. OSM
administrative units are not comparable internationally --- ``Riyadh,
Saudi Arabia'' returns a $375{,}478$\,km$^2$ province, ``Perth'' a
$6$\,km$^2$ council area, ``Alice Springs'' $1$\,km$^2$ --- so every
city, US and non-US alike, is defined as a $25$\,km-radius disc
($1{,}960$\,km$^2$, within the US sample's own range), with US cities
included under the identical definition as a bridge. Each city is sampled
in its own warm season.

Inversion is climatic. Nine cities outside the US show negative ceilings
(Table~\ref{tab:arid}), spanning Arabia, the Levant, North Africa, the
Atacama, Mexico and central Australia. Arid and humid/temperate ceilings
differ sharply (median $-0.032$ vs.\ $+0.615$, Mann--Whitney
$p{=}2.4{\times}10^{-4}$), while arid US and arid non-US cities do
\emph{not} differ detectably ($p{=}0.075$) --- the comparison an all-US
sample cannot make. The gap heuristic replicates on this independent
sample ($\rho{=}{+}0.92$ univariate, $+0.57$ multivariate), as does the
aridity link (mean LST vs.\ ceiling, $\rho{=}{-}0.69$).

\textbf{A robustness caveat we report rather than omit.} Inversion in the
American Southwest is sensitive to how the city is delimited. Of the five
US cities that invert under OSM city limits, four (Phoenix, Tucson,
Albuquerque, Boise) do \emph{not} invert under the disc definition; only
Billings does. The hyper-arid non-US cities invert decisively under either
(Antofagasta gap $-3.38$\,$^{\circ}$C, Alice Springs $-3.11$\,$^{\circ}$C), so the
mechanism is real, but the US examples sit near the sign boundary and
should not be read as definition-independent. Palm Springs is the
converse control: no inversion ($+0.255$) and the sample's largest
positive gap ($+6.34$\,$^{\circ}$C), in the same desert as Yuma, because its
landscape is irrigated --- the mechanism behaving as predicted.

Because the ceiling depends only on input data, candidate channels can be
screened without running the simulator at all. Adding MODIS NDVI is
predictable across the 49-city sample: gain correlates with ceiling
height ($\rho{=}{-}0.55$, $p{=}4.9\times10^{-5}$) and aridity
($\rho{=}{+}0.50$, $p{=}2.7\times10^{-4}$), with 27 of 49 cities robustly
positive (median gain $+0.06$). Phoenix --- the irrigated-desert city
motivating the channel --- ranks 4th of 49 ($+0.22$), behind Bakersfield,
Spokane, and Boise; this is a general, aridity-driven screening rule
rather than a single favorable city's result.

\subsection{Penman-Monteith Ablation}
Replacing linear ET with canonical Penman-Monteith reduces test $r$ from
$+0.755$ to $+0.416$. LAI saturation and VPD decoupling --- the two
mechanisms we isolated --- explain only 21\% of this regression combined;
stomatal resistance $r_s$ recovers 61\% at its optimum
($r_s{=}3200$\,s/m, $\sim32\times$ a reference grass value), but even
there PM ($r{=}{+}0.686$) underperforms simply \emph{removing} the ET
term ($r{=}{+}0.758$). No formulation we tested,
including one tuned to a physically indefensible parameter specifically
to maximize its own score, recovers what deleting the term already
achieves. This reframes the question from ``which ET formulation is
best'' to ``does this linear ET term do useful work at 500\,m at all'' ---
it evidently does not.

\section{Discussion}
\label{sec:disc}

\textbf{Diffusion trade-off.} No single Fickian diffusivity $D$ matches
both MODIS spatial autocorrelation and literature cooling-decay scales:
the best $r$ achievable inside the $100$--$300$\,m band is $+0.67$
($D{=}0.20$), a cost of $-0.19$ in $r$ against the unconstrained optimum
($r{=}{+}0.85$ at $D{\approx}12.6$) --- roughly twice the cost a naive,
uncalibrated estimate would suggest.

\textbf{Advection is not a drop-in fix.} Replacing diffusion with
ERA5 wind-driven advection destroys skill ($r{=}{-}0.11$ at 25\%
strength): a spatially uniform wind translates the whole surface pattern
downwind rather than the air layer that would physically transport heat;
a coupled air layer, not a one-line substitution, is required
\cite{mitraka2025prototypes}.

\textbf{Class separability as an a priori screen.} Because the ceiling
depends only on input data, it distinguishes cities where target-city
calibration can help (Denver) from those where it cannot (Phoenix)
\emph{before} any calibration is attempted --- a practical use beyond
simply explaining low scores after the fact.

\textbf{Limitations.} A $\sim$7\,$^{\circ}$C absolute bias (simplified longwave
scheme) couples weakly to skill ($\Delta r{=}0.02$ across a 3\,$^{\circ}$C bias
sweep); the 500\,m grid smooths highway-corridor hot streaks; and
per-city ceilings should be read as multi-day estimates, since single-day
values can be actively unrepresentative (Section~\ref{sec:results}).

\section{Conclusion}
\textsc{austin-twin} achieves $r{=}{+}0.754$ on held-out MODIS days using
only zero-authentication open data. Three contributions transfer beyond
this simulator: (1) a \emph{multivariate data-ceiling diagnostic} --- the
univariate alternative is unsound in arid cities --- validated across 49
cities and shown actionable via a predictable NDVI channel-screening
rule; (2) a \emph{diffusion--fidelity trade-off}, roughly twice as steep
as a naive estimate; and (3) an \emph{evapotranspiration ablation}
showing a standard linear ET term contributes nothing to spatial-pattern
skill at 500\,m, regardless of formulation. Code and validation scripts:
\url{https://github.com/abhi-is-cool/austin-twin}.

\section*{Acknowledgment}
MODIS data via Microsoft Planetary Computer; ESA WorldCover via Copernicus/ESA;
ERA5 via Google ARCO-ERA5 on GCS; OSM data copyright OpenStreetMap contributors (ODbL).

\bibliographystyle{IEEEtran}
% SPCONF 4/4 -- replace the line above with:
%   \bibliographystyle{IEEEbib}
\begin{thebibliography}{14}

\bibitem{bhamjee2024detection}
M.~Bhamjee et al., ``Detection and characterization of urban heat islands
with machine learning,'' in \textit{Proc. IGARSS 2024}, pp.~1693--1699.

\bibitem{bueno2013urban}
B.~Bueno, L.~Norford, J.~Hidalgo, and G.~Pigeon, ``The urban weather
generator,'' \textit{J. Building Perform. Simul.}, vol.~6, no.~4,
pp.~269--281, 2013.

\bibitem{hersbach2020era5}
H.~Hersbach et al., ``The ERA5 global reanalysis,''
\textit{Quart. J. Royal Meteorol. Soc.}, vol.~146, pp.~1999--2049, 2020.

\bibitem{hoang2025geospatial}
N.~D.~Hoang and Q.~L.~Nguyen, ``Geospatial analysis and machine learning
framework for urban heat island intensity prediction,''
\textit{Sustainability}, vol.~17, p.~4287, 2025.

\bibitem{kusaka2001simple}
H.~Kusaka et al., ``A simple single-layer urban canopy model,''
\textit{Boundary-Layer Meteorol.}, vol.~101, pp.~329--358, 2001.

\bibitem{lyu2022integrated}
F.~Lyu et al., ``An integrated CyberGIS and machine learning framework for
fine-scale UHI prediction,'' \textit{Urban Informatics}, vol.~1, p.~6, 2022.

\bibitem{masson2000physically}
V.~Masson, ``A physically-based scheme for the urban energy budget,''
\textit{Boundary-Layer Meteorol.}, vol.~94, pp.~357--397, 2000.

\bibitem{microsoft_pc}
Microsoft, ``Planetary Computer STAC API,'' 2023. [Online].
Available: \url{https://planetarycomputer.microsoft.com}

\bibitem{mitraka2025prototypes}
Z.~Mitraka et al., ``Prototypes for enhanced urban heat monitoring,''
\textit{Discover Cities}, vol.~2, p.~14, 2025.

\bibitem{oke1982energetic}
T.~R.~Oke, ``The energetic basis of the urban heat island,''
\textit{Quart. J. Royal Meteorol. Soc.}, vol.~108, pp.~1--24, 1982.

\bibitem{storn1997differential}
R.~Storn and K.~Price, ``Differential evolution -- a simple and efficient
heuristic,'' \textit{J. Global Optim.}, vol.~11, pp.~341--359, 1997.

\bibitem{voogt2003thermal}
J.~A.~Voogt and T.~R.~Oke, ``Thermal remote sensing of urban climates,''
\textit{Remote Sens. Environ.}, vol.~86, pp.~370--384, 2003.

\bibitem{zanaga2022esa}
D.~Zanaga et al., ``ESA WorldCover 10\,m 2021 v200,'' \textit{Zenodo}, 2022.

\bibitem{ziter2019scale}
C.~D.~Ziter et al., ``Scale-dependent interactions between tree canopy cover
and impervious surfaces,'' \textit{Proc. Natl. Acad. Sci.}, vol.~116,
pp.~7575--7580, 2019.

\end{thebibliography}

\end{document}
